\QHeading{第1题　在$\square$里填上合适的数。}
\par\noindent


\Equation{(20+4)\times 25=\textcolor{answerred}{20}\times 25+\textcolor{answerred}{4}\times 25}
\par

\Equation{16\times (30+5)=16\times\textcolor{answerred}{30}+16\times\textcolor{answerred}{5}}
\Equation{47\times 15+53\times 15=(\textcolor{answerred}{47}+\textcolor{answerred}{53})\times 15}
\Equation{78\times 24+22\times 24=\textcolor{answerred}{24}\times (\textcolor{answerred}{78}+\textcolor{answerred}{22})}
\Equation{8\times (125+50)=\textcolor{answerred}{8}\times\textcolor{answerred}{125}+\textcolor{answerred}{8}\times\textcolor{answerred}{50}}
\Equation{63\times 99+63=63\times (\textcolor{answerred}{99}+\textcolor{answerred}{1})}
\par\noindent

\QHeading{第2题　用简便方法计算下面各题。}

\CalcBlock[方法一：简便计算]{1}{\begin{aligned}[t]&\phantom{={}}\color{black}42\times (100+5)\\&\color{answerred}=42\times100+42\times5\\&\color{answerred}=4200+210\\&\color{answerred}=4410\end{aligned}}
\par
\MethodBlock{方法二：先处理括号或乘除，再凑整后补偿}{把接近整十、整百的因数凑整，再补回或减去多算的部分。}{\begin{aligned}[t]&\phantom{={}}42\times (100+5)\\
&=42\times 105\\
&=105\times (40+2)\\
&=105\times 40+105\times 2\\
&=4200+210\\
&=4410\end{aligned}}


\CalcBlock[方法一：简便计算]{2}{\begin{aligned}[t]&\phantom{={}}\color{black}99\times 101-99\\&\color{answerred}=99\times(101-1)\\&\color{answerred}=99\times100\\&\color{answerred}=9900\end{aligned}}
\MethodBlock{方法二：先处理括号或乘除，再多减再加回}{}{\begin{aligned}[t]&\phantom{={}}99\times 101-99\\
&=9999-99\\
&=9999-100+1\\
&=9899+1\\
&=9900\end{aligned}}

\CalcBlock[方法一：简便计算]{3}{\begin{aligned}[t]&\phantom{={}}\color{black}34\times 127-27\times 34\\&\color{answerred}=34\times127-34\times27\\&\color{answerred}=34\times(127-27)\\&\color{answerred}=34\times100\\&\color{answerred}=3400\end{aligned}}
\MethodBlock{方法二：先处理括号或乘除，再多减再加回}{}{\begin{aligned}[t]&\phantom{={}}34\times 127-27\times 34\\
&=4318-918\\
&=4318-920+2\\
&=3398+2\\
&=3400\end{aligned}}

\CalcBlock[方法一：简便计算]{4}{\begin{aligned}[t]&\phantom{={}}\color{black}37\times 99\\&\color{answerred}=37\times(100-1)\\&\color{answerred}=37\times100-37\times1\\&\color{answerred}=3700-37\\&\color{answerred}=3663\end{aligned}}
\MethodBlock{方法二：按数位拆分计算}{把一个因数拆成整十、整百数与余下部分之和。}{\begin{aligned}[t]&\phantom{={}}37\times 99\\
&=37\times (90+9)\\
&=37\times 90+37\times 9\\
&=3330+333\\
&=3663\end{aligned}}

\CalcBlock[方法一：简便计算]{5}{\begin{aligned}[t]&\phantom{={}}\color{black}(49\times 25)\times 4\\&\color{answerred}=49\times(25\times4)\\&\color{answerred}=49\times100\\&\color{answerred}=4900\end{aligned}}
\MethodBlock{方法二：按运算顺序直接计算（对照）}{本法只作结果核对；本题要求简便计算，前面的简便解法仍保留。}{\begin{aligned}[t]&\phantom{={}}(49\times 25)\times 4\\
&=1225\times 4\\
&=4900\end{aligned}}

\CalcBlock[方法一：简便计算]{6}{\begin{aligned}[t]&\phantom{={}}\color{black}65\times 102\\&\color{answerred}=65\times(100+2)\\&\color{answerred}=65\times100+65\times2\\&\color{answerred}=6500+130\\&\color{answerred}=6630\end{aligned}}
\MethodBlock{方法二：按数位拆分计算}{把一个因数拆成整十、整百数与余下部分之和。}{\begin{aligned}[t]&\phantom{={}}65\times 102\\
&=102\times (60+5)\\
&=102\times 60+102\times 5\\
&=6120+510\\
&=6630\end{aligned}}

\par\noindent

\QHeading{第3题　解决问题。}

\Problem{① 甲、乙两个工程队合修一条公路，甲队每天修45米，乙队每天修35米，两队同时开工，12天修完。这条公路长多少米?}
\par


\MethodTitle{方法一}
\Answer{\Formula{\begin{aligned}[t]&\phantom{={}}(45+35)\times12\\&=80\times12\\&=960\text{（米）}\end{aligned}}}
\Answer{答：这条公路长960米。}

\MethodTitle{方法二：分别计算两部分，再求和}
\Answer{不先合并每份的数量，而是分别求出两部分的总数，再相加。}
\Answer{\Formula{\begin{aligned}[t]&\phantom{={}}45\times 12+35\times 12\\
&=540+420\\
&=960\text{（米）}\end{aligned}}}
\Answer{答：这条公路长960米。}
\Problem{② 超市搞促销，每千克苹果8元，每千克香蕉12元，妈妈买了5千克苹果和5千克香蕉，一共花了多少钱?（用两种方法解答。）}
\Answer{方法一：}
\Answer{\Formula{\begin{aligned}[t]&\phantom{={}}8\times5+12\times5\\&=40+60\\&=100\text{（元）}\end{aligned}}}
\Answer{方法二：}
\Answer{\Formula{\begin{aligned}[t]&\phantom{={}}(8+12)\times5\\&=20\times5\\&=100\text{（元）}\end{aligned}}}
\Answer{答：一共花了100元。}
